\section{Límites y futuras líneas de investigación}
\label{sec:mp-limites}

% Fuente única de cifras: notebooks/STRATA_marco_practico.ipynb (panel canónico de 10) y sus JSON en
% outputs/experiments/. Toda cifra trazada a su JSON en el caption de su tabla o en comentario LaTeX.

Un rescate probado necesita sus límites dichos con la misma claridad. Reconocer dónde acaba el alcance vuelve defendible lo demostrado.

\subsection{Qué no sobrevive a un test}
\label{sec:mp-lim-tests}

Cinco resultados del capítulo viven por debajo del listón de la significancia. El primero es la ventaja en accuracy sobre la referencia trivial: en SPY, AutoML acierta el $57{,}4\,\%$ frente al $56{,}6\,\%$ de ZeroR, pero el McNemar pareado da $p = 0{,}902$ y en doscientos cincuenta días no hay potencia para separarlos. En el resto del panel ninguna derivada bate a la clase mayoritaria con significancia, y la estrategia ganadora cambia de un activo a otro. La victoria es real en el número, nominal en el test y leída a posteriori. El segundo es la ley del leverage: el rescate escala con el \emph{leverage effect} del activo, con correlación de Pearson $r = -0{,}56$ ($p = 0{,}093$) y de Spearman $\rho = -0{,}59$ ($p = 0{,}074$) sobre los diez activos. Queda por debajo de $0{,}10$, el nivel que el proyecto fijó por adelantado, pero no alcanza el $0{,}05$ convencional ni resiste un \emph{leave-one-out}; con diez puntos la correlación es fuerte y lo que falta es potencia. Es una tendencia, no una ley cerrada.

El tercero pesa sobre la regla. El rescate de riesgo agregado, medido por el $\Delta$Sharpe \emph{pooled} de los diez activos frente al agente, excluye el cero para los tres supervisores ($+0{,}60$ M8, $+1{,}12$ M10, $+1{,}08$ AutoML). Bajo una cota de Bonferroni con $m = 3$ pruebas (López de Prado 2018 \cite{lopezdeprado2018}), M10 y AutoML siguen por encima del cero, pero M8 cae a $-0{,}047$: la protección de la regla no resiste el ajuste por número de pruebas. El cuarto es una falsación pre-registrada. En SPY en solitario, restringiendo al régimen bajista ($n = 50$ días), M8 invierte su signo y su $\Delta$Sharpe frente al agente cae a $-1{,}49$. Anotamos esta condición de fracaso antes de mirar el resultado, y la agregación del panel la resuelve: el rescate de riesgo es un fenómeno de conjunto, no de un activo aislado. El quinto afecta al clustering: los cuatro métodos recuperan los mismos tres grupos con un índice de Rand de $1{,}0$, pero sobre diez activos, así que es exploratorio y no permite predecir qué modelo ganará en un activo nuevo. La Tabla~\ref{tab:mp-limites} reúne los cinco con su test y su veredicto.

\begin{table}[htbp]
\centering
\caption{Resumen de los resultados que no alcanzan significancia plena, con su contraste y su veredicto. Cada fila enlaza una afirmación del capítulo con el test que la mide y con el JSON del que sale la cifra. La columna de veredicto distingue lo nominal de lo falsado y de lo exploratorio. Fuentes: \texttt{automl\_runs/panel\_mm25\_inclGBM-XGB-SE\_AUC\_emb1\_N0-150\_step21\_kfold\_seed42.json}, \texttt{leverage\_law\_panel10.json}, \texttt{bullbear\_confirmatory.json} (\texttt{POOLED10}), \texttt{cluster\_panel10.json}.}
\label{tab:mp-limites}
\begin{tabular}{p{4.4cm}p{4.6cm}p{4.0cm}}
\toprule
Afirmación & Test y cifra & Veredicto \\
\midrule
Una derivada bate a la trivial en accuracy & McNemar AutoML vs ZeroR sobre SPY, $p = 0{,}902$ & Nominal ($n \approx 250$, a posteriori) \\
\addlinespace
El rescate en accuracy escala con el \emph{leverage} & Pearson $r = -0{,}56$, $p = 0{,}093$; Spearman $\rho = -0{,}59$, $p = 0{,}074$ (diez activos) & Tendencia al $\alpha = 0{,}10$; no $p < 0{,}05$ ni \emph{leave-one-out} \\
\addlinespace
La regla M8 rescata el riesgo agregado & $\Delta$Sharpe \emph{pooled} $+0{,}60$, cota inferior de Bonferroni ($m = 3$) $= -0{,}047$ & No sobrevive a la corrección múltiple \\
\addlinespace
La regla M8 rescata el riesgo en SPY-bajista & $\Delta$Sharpe $= -1{,}49$ ($n = 50$) & Falsación pre-registrada, resuelta por el panel \\
\addlinespace
La naturaleza define grupos de activos & Consenso de cuatro métodos, Rand $= 1{,}0$ (diez activos) & Exploratorio, no predictivo por activo \\
\bottomrule
\end{tabular}
\end{table}

Dos cautelas atraviesan estos límites. La ventana corta: con $n \approx 250$ días la significancia plena frente a la trivial está fuera de alcance por potencia, no por ausencia de señal. Y el \emph{pooled}: apilar los días de los diez activos los trata como independientes, y no lo son, porque la correlación cruzada entre índices como SPY y QQQ supera $0{,}9$ y el tamaño efectivo de muestra es menor que el nominal. Lo controlamos con un bootstrap estacionario de bloque medio $\sqrt{N}$ (Politis y Romano 1994 \cite{politis1994}), y aun así el $\Delta$Sharpe de M10 y AutoML excluye el cero con holgura; la regla vive en el filo.

\subsection{Desplegabilidad}
\label{sec:mp-lim-despliegue}

Los tres supervisores no cuestan lo mismo de poner en marcha. La regla M8 opera desde el primer día: sus umbrales se fijan en la calibración de 2000 a 2024 y no necesita reentreno, sin ventana de arranque ni modelo que mantener. Los aprendices piden más. M10 y AutoML entrenan en un \emph{walk-forward} de origen rodante y necesitan unos ciento cincuenta días de arranque antes de su primera decisión, y luego reentrenan de forma periódica. AutoML añade una carga extra: en cada reentreno H2O vuelve a buscar entre GBM, XGBoost y su apilamiento, y el modelo líder cambia de una iteración a otra, lo que obliga a auditar un objeto distinto cada vez en lugar de vigilar una pieza estable.

El balance es claro. La regla es la más sencilla de desplegar, lista para producción sin arranque ni reentreno; los aprendices dan más accuracy a cambio de más mantenimiento y de una pieza que cambia bajo los pies. Quien priorice robustez operativa se queda con la regla; quien priorice precisión direccional paga el coste de gobernar un aprendiz que se reentrena. Ninguno de estos balances incorpora todavía costes de transacción, que un despliegue real tendría que medir.

\subsection{Alcance y líneas futuras}
\label{sec:mp-lim-futuro}

El resultado que sobrevive a un test es el rescate del agente: STRATA es una capa de supervisión y de control de riesgo sobre un agente que pierde, y se mide por cuánto lo mejora. Superar al mercado pasivo nunca fue su objetivo, y no lo logra con significancia, porque en la ventana de $n \approx 250$ días no hay potencia para separar a las derivadas de la referencia pasiva. En punto sí asoma un valor que merece nombrarse: en cuatro de los diez activos una derivada de STRATA supera al pasivo en Sharpe, y lo hace justamente donde el \emph{leverage} es débil o se invierte. En SMCI, MARA y UNG el comprar y mantener pierde o queda plano, y la derivada, al ir defensiva en el momento correcto, saca un Sharpe positivo claro. No es exposición pasiva capturada por casualidad: es compatible con valor direccional, pendiente de un contraste pre-registrado. Sigue siendo nominal y leído a posteriori, y por eso no asciende a conclusión, pero señala que el camino existe.

Ese camino organiza las líneas futuras. La vía más directa es repetir el estudio sobre una muestra mayor, en activos y en días, para que la significancia frente a la trivial deje de estar limitada por potencia. Otra surge de la complementariedad por régimen: M10 rescata más en mercado alcista y AutoML en bajista, con un contraste de diferencia en diferencias significativo sobre el panel, lo que sugiere un ensemble enrutado por régimen. Cabe también abrir el banco de pruebas a otros agentes y a otros modelos de lenguaje, para separar lo propio de STRATA de lo propio de este agente concreto, y llevar el sistema a un despliegue en vivo. La más ambiciosa queda fuera del cuerpo de esta memoria: desarrollar a partir de STRATA una estrategia que genere alfa robusta, una rentabilidad absoluta que bata al mercado con un contraste pre-registrado. Las victorias direccionales en SMCI, MARA y UNG son su punto de partida, no su conclusión.

STRATA rescata al agente en accuracy y en riesgo con significancia, el aprendiz redescubre sus señales y mejora modestamente a la regla, y la naturaleza del activo explica qué canal rescata. El listón sigue siendo el agente, y eso es lo que el trabajo demuestra.
